\section{Conclusion}

We studied a defect-aware reliability layer for assistive VQA using real answerability and image-quality labels from VizWiz. Lightweight heads on frozen visual backbones can predict answerability and diagnose visual defects, and these predictions support actionable retake guidance. At the same time, defect-aware risk scores do not outperform global VQA confidence for selective prediction. The resulting lesson is practical: confidence should remain the primary risk-ranking signal, while defect diagnosis should be used to explain refusal and guide recovery.
